\section{Methodology}
\label{sec:method}

\subsection{Benchmark harness}
\label{sec:harness}

\sysname is an asynchronous Python client that drives each engine through its own
OpenAI-compatible \texttt{/v1/completions} endpoint with streaming enabled, so every engine is
measured through the interface applications use, including its HTTP frontend and tokenizer.
For each request the harness records the dispatch time, the arrival time of every streamed
chunk, the number of output tokens, and the final status. In parallel it polls each engine's
native metrics endpoint (Prometheus \texttt{/metrics} for \vllm and \sglang, a JSON statistics
endpoint for \trt) for KV-cache occupancy, running and waiting sequence counts, preemption
counters, and prefix-cache hit rates.

\paragraph{Load generation.}
Arrival times are generated in advance and requests are released open-loop, independent of
completions~\cite{schroeder2006open}. We use four Poisson processes (2, 4, 8, and 16 requests/s)
and one bursty process (10 requests 50\,ms apart, every 5\,s). A client-side semaphore caps the
number of in-flight requests (\emph{client concurrency}, 16 or 256). For multi-turn
conversations, a turn is released only after the previous turn of the same conversation
completes, as a real user would; conversations are interleaved open-loop.

\paragraph{Isolation.}
Each rate configuration runs against a freshly launched server, so caches and counters start
empty. \todo{state warm-up policy and inter-run cool-down} Requests that receive no first token
within 300\,s, or do not finish within 600\,s, are recorded as timeouts.

\subsection{Engines}
\label{sec:engines}

\Cref{tab:engines} lists the engines and versions. All engines serve fp16 weights on one GPU
(tensor parallelism 1) with a maximum context of 8{,}192 tokens. The scheduler's sequence
limit (\texttt{max\_num\_seqs}) is set to 64 or 256 and mapped to each engine's equivalent
setting. \todo{list the per-engine flags: KV memory fraction, maximum batched tokens,
maximum running requests, maximum batch size, and the capacity-scheduler policy}
Two engine-specific adjustments were necessary. First, \trt treats the chunk size as a hard
prompt-length limit unless chunked prefill is enabled explicitly, so the chunked-prefill
strategy sets this flag. Second, \trt's n-gram speculation is incompatible with its overlap
scheduler, which we therefore disable whenever n-gram speculation is on; \sglang disables its
equivalent optimisation automatically and caps running requests at 48 when speculation is on.

\begin{table}[t]
  \caption{Engines and platforms.}
  \label{tab:engines}
  \small
  \begin{tabular}{@{}lll@{}}
    \toprule
    & L40S (\rqref{1}--\rqref{3}) & GB10 (\rqref{4}) \\
    \midrule
    \vllm (V1) & 0.11.0 & 0.32.0 \\
    \sglang & \verify{version} & 0.5.22 \\
    \trt (PyTorch) & 1.0.0 & 1.2.1 \\
    \midrule
    Memory & 44.4\,GiB & 121\,GiB unified \\
    KV memory fraction & 0.90 \verify{} & 0.80 \\
    \bottomrule
  \end{tabular}
\end{table}

\subsection{Models}
\label{sec:models}

We select three instruction-tuned models whose per-token KV footprint spans a wide range
(\cref{tab:models}). Qwen1.5-14B-Chat~\cite{bai2023qwen} is the only one without any KV-head
reduction ($H_{kv}=40$), so it needs $5\times$ the cache per token of Mistral-NeMo-12B and
$6.25\times$ that of Llama-3.1-8B~\cite{dubey2024llama3}. We considered distilled reasoning
models but rejected them after confirming from their configuration files that they use GQA
backbones and would not provide the architectural contrast.

\begin{table}[t]
  \caption{Models. KV bytes per token at fp16.}
  \label{tab:models}
  \small
  \begin{tabular}{@{}llrrr@{}}
    \toprule
    Model & Attn. & $H_{kv}$ & Layers & KV B/token \\
    \midrule
    Llama-3.1-8B-Instruct & GQA & 8 & 32 & 131{,}072 \\
    Mistral-NeMo-12B-Instr.~\cite{mistralnemo} & GQA & 8 & 40 & 163{,}840 \\
    Qwen1.5-14B-Chat & MHA & 40 & 40 & 819{,}200 \\
    \bottomrule
  \end{tabular}
\end{table}

\paragraph{Prompt formatting.}
The harness sends plain instruction text to the completions endpoint. Qwen1.5-14B-Chat emits
its end-of-turn token immediately on unformatted prompts, so for this model the harness
applies the model's chat template at send time. \todo{decide whether to template all models;
see \cref{sec:threats} on output length and zero-token replies}

\subsection{Workloads}
\label{sec:workloads}

All workloads are built from public corpora (\cref{tab:workloads}).
\texttt{mixed\_real} mixes four task types in equal shares: chat prompts from
ShareGPT~\cite{sharegpt}, code generation from MBPP~\cite{austin2021mbpp}, reasoning from
GSM8K~\cite{cobbe2021gsm8k}, and summarisation from CNN/DailyMail~\cite{hermann2015cnn}. It
combines short prompts with very long requested outputs (reasoning) and long prompts with short
outputs (summarisation), so a few long-running requests can occupy scheduler slots and KV memory
for minutes. \texttt{chat\_multiturn} replays ShareGPT conversations of two to six turns; each
turn resends the conversation history, which makes prefix reuse possible. \texttt{chat\_real\_long}
uses single-turn ShareGPT prompts of at least 500 tokens to stress prefill.

\begin{table}[t]
  \caption{Workloads (per rate configuration; token counts are medians).}
  \label{tab:workloads}
  \footnotesize
  \setlength{\tabcolsep}{4pt}
  \begin{tabular}{@{}llrrr@{}}
    \toprule
    Workload & Task (source) & Req. & Prompt & Max out \\
    \midrule
    \multirow{4}{*}{\texttt{mixed\_real}}
      & chat (ShareGPT) & 125 & 20 & 239 \\
      & code (MBPP) & 125 & 53 & 45 \\
      & reasoning (GSM8K) & 125 & 85 & 6{,}810 \\
      & summ.\ (CNN/DM) & 125 & 791 & 64 \\
    \texttt{chat\_multiturn} & 2--6 turns (ShareGPT) & 280--600 & 454 & 310 \\
    \texttt{chat\_real\_long} & $\geq$500 tok.\ (ShareGPT) & 300 & 938 & \verify{} \\
    \bottomrule
  \end{tabular}
\end{table}

\subsection{Metrics}
\label{sec:metrics}

\begin{description}
  \item[TTFT] Time to first token. For open-loop \texttt{mixed\_real} we measure it from the
  request's scheduled \emph{arrival}, i.e., including any time spent queued at the client;
  measuring from dispatch hides that queueing (\cref{sec:discussion}). For
  \texttt{chat\_multiturn}, where later turns are deliberately held until the previous turn
  ends, we measure from dispatch.
  \item[TBT] Time between consecutive streamed chunks, summarised per request by its 95th
  percentile and maximum (the longest stall).
  \item[TPOT] Time per output token, $(\mathrm{E2E}-\mathrm{TTFT})/(n-1)$ for $n$ output tokens.
  Unlike TBT it is unaffected by how many tokens a chunk carries, which matters under
  speculative decoding.
  \item[E2E] End-to-end latency, from arrival to the last token.
  \item[Throughput] Output tokens per second over a rate configuration.
  \item[Outcome] Finished, timed out, or \emph{zero-token reply} (the server ends the stream
  without producing a token).
  \item[Goodput] The share of \emph{all} requests that finish while meeting both a TTFT and a
  TPOT target. We use a \emph{strict} SLO (TTFT $\le$ 2\,s, TPOT $\le$ 200\,ms) and a
  \emph{relaxed} SLO (TTFT $\le$ 4\,s, TPOT $\le$ 500\,ms).
\end{description}

\paragraph{Censoring.}
Discarding timed-out requests would make an overloaded engine look fast. We therefore count a
request with no first token by 300\,s as TTFT $=$ 300\,s and E2E $\ge$ 300\,s, and a request
that started but did not finish by 600\,s as E2E $=$ 600\,s. Zero-token replies are excluded
from latency statistics and count as goodput misses.

\paragraph{Scheduler-internal metrics.}
Engine-reported metrics are used only to explain user-facing results, because their semantics
differ between engines: \vllm reports a cumulative prefix-cache hit rate and a per-server
preemption counter; \sglang reports the hit rate of the most recent batch and, from version 0.5,
a retraction counter; \trt's statistics endpoint is read once after each run in our harness.

\subsection{Experimental design}
\label{sec:design}

\Cref{tab:design} summarises the configurations per research question. Every configuration is
run with all three engines and all five arrival processes.

\begin{table}[t]
  \caption{Experimental design.}
  \label{tab:design}
  \small
  \begin{tabular}{@{}lp{6.1cm}@{}}
    \toprule
    RQ & Configurations \\
    \midrule
    \rqref{1} & 3 models; \texttt{mixed\_real}, \texttt{chat\_multiturn}; client concurrency 16, 256
    (256 only for Qwen1.5); \texttt{max\_num\_seqs} 64, 256; baseline strategy \\
    \rqref{2} & Mistral-NeMo; \texttt{chat\_multiturn}, \texttt{chat\_real\_long}; prefix caching
    on/off $\times$ chunked prefill off/512 tokens; concurrency 256 \\
    \rqref{3} & Mistral-NeMo; \texttt{mixed\_real}, \texttt{chat\_multiturn}; baseline vs.\ n-gram
    ($k{=}3$); \texttt{max\_num\_seqs} 32, 64; concurrency 256 \\
    \rqref{4} & \rqref{1} protocol on GB10 at concurrency 256 \\
    \bottomrule
  \end{tabular}
\end{table}

In total the study comprises \todo{count} rate configurations and \todo{count} requests on the
L40S and 180 rate configurations and 72{,}000 requests on the GB10. Each configuration is run
once.
